\documentclass[a4paper]{article}
\usepackage[pdftex]{graphicx}
\usepackage[utf8]{inputenc}
\usepackage{enumerate}
\usepackage{siunitx}
\sisetup{locale=DE} 
\usepackage{href-ul}
\hypersetup{
	colorlinks=true,
	linkcolor=blue,
	urlcolor=blue}
\usepackage{icomma}
\usepackage{amssymb}
\usepackage{geometry}
\geometry{a4paper, top=15mm, left=15mm, right=15mm, bottom=15mm,
	headsep=10mm, footskip=12mm}

%Source: img_0096

% Start the document
\begin{document}
%Titel
{\bf Schulaufgabe aus der Physik }\\
\begin{enumerate}[1.]

%Aufgabe 1
{\bf \item In einem Experiment wird die Abhängigkeit der Stromstärke $I$ von der Spannung $U$ für einen Eisen-- und einen Konstantandraht untersucht.}

\begin{enumerate}[a)]
\item Fertigen Sie eine Versuchsskizze (=Schaltskizze) an.
\item In einem Versuch ergeben sich folgende Messwerte:

\hspace{15 pt}

\begin{tabular}{c|c|c|c|c|c|c|c|c|c}
                             & $U$ in $\SI{}{\volt}$ & 0 & 0,30 & 0,60 & 0,90 & 1,20 & 1,50 & 1,80 & 2,10 \\
\hline
Eisendraht            & $I$ in $\SI{}{\ampere}$ & 0 & 0,48 & 0,93 & 1,38 & 1,70 & 1,98 & 2,19 & 2,31 \\
Konstantandraht  & $I$ in $\SI{}{\ampere}$ & 0 & 0,09 & 0,18 & 0,27 & 0,36 & 0,45 & 0,54 & 0,63 
\end{tabular}

\hspace{15 pt}

 Werten Sie die Messreihe graphisch und rechnerisch aus.
\item Interpretieren Sie den Verlauf der beiden Leiterkennlinien.
\item Erklären Sie mithilfe des Teilchenmodells den Verlauf der Kennlinie des verwendeten Eisendrahtes
\end{enumerate}

%Aufgabe 2
{\bf \item Ein Strommessgerät mit dem Innenwiderstand $R_i=\SI{70}{\ohm}$ kann Stromstärken bis 
$\SI{3,0}{\milli\ampere}$ messen.} Der Messbereich soll auf $\SI{30}{\milli\ampere}$ erweitert werden. Geben Sie an, wie ein zusätzlicher Widerstand geschaltet werden muss, und berechnen Sie dessen Wert.

%Aufgabe 3
{\bf \item Im Anlasser eines Automotors fließt beim Betrieb ein Strom von $I=\SI{120}{\ampere}$.} Die Starterbatterie hat einen Innenwiderstand von $R_i=\SI{0,030}{\ohm}$ und eine Quellenspannung von $U=\SI{12}{\volt}$\\
Berechnen Sie die Betriebsspannung während des Anlassvorgangs sowie den Widerstand  des Anlassers und dessen elektrische Leistung.

%Aufgabe 4
{\bf \item Zwei Heizspiralen einer elektrischen Kochplatte haben jeweils einen Widerstand von 
$R=\SI{50}{\ohm}$.} Mit diesen Heizspiralen lassen sich durch geeignete Schaltungen drei Stufen unterschiedlicher Heizleistung erreichen.

\begin{enumerate}[a)]
\item Zeichnen Sie für die Stufe mit kleinster und für die Stufe mit größter Heizleistung jeweils ein Schaltbild.
\item Berechnen Sie für beide Stufen jeweils den Gesamtwiderstand der Schaltung sowie die elektrische Leistung. (Spannung $\SI{230}{\volt}$).
\end{enumerate}

\end{enumerate} 

\fbox{
	\begin{minipage}{0.5\textwidth}
		Zur Lösung bitte \href{https://www.okuyakl.de/phys/p10elL096/ll096.pdf}{hier klicken} oder den QR-Code scannen.\\
	Weitere Arbeitsblätter gibt es unter 
	
	\href{https://www.okuyakl.de}{www.okuyakl.de}
	\end{minipage}
	\hfill
	\begin{minipage}{0.4\textwidth}
		\includegraphics[width=1.5 cm]{../../viecher/zwe03}
		\includegraphics[width=3 cm]{qr096}
		\includegraphics[width=2 cm]{../../viecher/afanticon1}
		
	\end{minipage}}

\end{document}%Lösung---------------------------------------
